\documentclass[11pt]{article}
\usepackage{fontspec}
\usepackage[french]{babel}
\usepackage{amsmath}
\usepackage{amssymb}
\usepackage{amsthm}
\usepackage{mathrsfs}
\usepackage{tikz-cd}
\usepackage{enumitem}

\newcommand{\printedpage}[1]{\marginpar{\footnotesize p.~#1}}
\newcommand{\msp}[1]{\marginpar{\footnotesize [#1]}}
\newcommand{\M}{\mathcal{M}}
\newcommand{\cercle}[1]{\tikz[baseline=(X.base)]\node[draw,circle,inner sep=0.6pt](X){#1};}
\newcommand{\Filt}{\mathrm{Filt}}
\newcommand{\Mon}{\mathrm{Mon}}
\newcommand{\Ss}{\mathrm{Ss}}
\newcommand{\Fl}{\mathrm{Fl}}
\newcommand{\Ob}{\mathrm{Ob}}
\newcommand{\Cof}{\mathrm{Cof}}

\begin{document}

\printedpage{44}
% head: A. Grothendieck : Dérivateurs, Ch. XV, Édition des Universités de Montpellier II et Paris VII

\msp{page 80}
ensemble petit, à cause des hypothèses faites sur $\M$), les ordinaux $\alpha$ pour lesquels $A_\alpha$ existe satisfont $\mathrm{card}\,A_\alpha \leq \pi$. Donc il doit exister un $\alpha$ tel que $A_\alpha = B$. Mais alors $i_\alpha : A \longrightarrow A_\alpha$ est égal à $f$, donc $f \in \widetilde{C^{\pi}}$, \quad \textit{q.e.d.}

Il reste donc à prouver 6.4. Pour ceci, soit $I$ l'ensemble des sous-objets de $B$ qui sont dans $\Filt^{\pi}$, ordonné par inclusion. (Si $B$ est dans $\Filt^{\pi}$, $A$ aussi, donc $f$ est déjà dans $C^{\pi} \subset \widetilde{C^{\pi}}$, et il n'y a rien à prouver.) C'est un ensemble filtrant grand devant $\pi$ (et même $\pi$-\textit{adapté}), dont la $\varinjlim$ est $B$ (5.3, p. 46). Considérons les
\[
A_i \;=\; B_i \cap A \;,
\]
ils sont dans $\Filt^{\pi}$, comme sous-objets d'objets de $\Filt^{\pi}$ (condition $(d)$ sur la filtration cardinale). Ainsi, la flèche $f$ apparaît comme limite inductive filtrante des flèches
\[
f_i \;:\; A_i \longrightarrow B_i \;.
\]
Il est trivial que le système des $A_i$ est cofinal dans $\Ss_{\pi}(A)$, donc $A = \varinjlim A_i$, d'où $f = \varinjlim f_i$.

\msp{page 81}
D'autre part, je dis que le système inductif des $f_i$ dans $\underline{\Fl}(\M)$ est \textit{bien adapté à $\pi$.} On sait déjà que $I$ est bien adapté à $\pi$, le sup dans $I$ d'une partie filtrante $J$ de cardinal $\leq \pi$ étant le sup dans l'ensemble ordonné $\Ss(B)$ (lequel sup se trouve dans $\Ss_{\pi}(B)$). Il est essentiel, pour notre argument, de savoir que si $i_J = \sup J$, alors non seulement $B_{i_J} = \varinjlim\limits_{i\in J} B_i$, mais \marginpar{$\|$}% marginal double vertical bar printed in the left margin at this line
aussi $A_{i_J} = \varinjlim\limits_{i\in J} A_i$, de sorte que le système inductif des $f_i$ est $\pi$-adapté dans $\underline{\Fl}(\M)$. Il semble qu'il faille donc une condition d'exactitude sur $\M$ voisine de $(e)$, savoir :

\begin{itemize}
\item[\cercle{e$'$}] Les foncteurs $\varinjlim$ filtrantes commutent au changement de base.
\end{itemize}

(Il le suffit ici dans le cas très particulier d'un changement de base par un monomorphisme $A \longrightarrow B$, et d'une limite inductive filtrante de sous-objets de $B$ pour des morphismes d'inclusion comme transitions.) Notons que \cercle{e}, \cercle{e$'$} se trouvent généralement comme conséquence de la condition % [sic] Il le suffit -> Il suffit, ungrammatical as printed

\begin{itemize}
\item[(e$_0$)] Les $\varinjlim$ filtrantes dans $I$ sont exactes à gauche.
\end{itemize}

\msp{page 82}
Ceci posé, on voit donc, $C$ étant $\pi$-accessible et $I$ grand devant $\pi$, que l'ensemble $I_C$ des $i \in I$ tels que $f_i \in C$ est \textit{cofinal} dans $I$. Mais comme $f_i \in \Filt^{\pi}$, $f_i \in C$ équivaut à $f_i \in C^{\pi}$. Comme $A \longrightarrow B$ n'est pas un isomorphisme, il existe un $i_0$ tel que $B_{i_0} \not\subset A$, puis un $i \geq i_0$ tel que $f_i \in C^{\pi}$, et aussi $B_i \not\subset A$ (puisque $B_i \supset B_{i_0}$). Ainsi on a prouvé ceci :

\begin{itemize}
\item[($*$)] Sous les hypothèses de (6.4) et moyennant l'hypothèse \cercle{e$'$} sur $\M$, il existe un sous-objet $B'$ de $B$ satisfaisant les conditions suivantes.
\begin{itemize}
\item[\textit{a})] $B' \in \Filt^{\pi}$.
\item[\textit{b})] $B' \not\subset A$.
\item[\textit{c})] Posant $A' = B' \cap A$, l'inclusion
\[
i' \;:\; A' \;\longrightarrow\; B'
\]
est dans $C^{\pi}$.
\end{itemize}

\printedpage{45}
% head: A. Grothendieck : Dérivateurs, Ch. XV, Édition des Universités de Montpellier II et Paris VII

De plus, l'ensemble des $J \subset \Ss_{\pi}(B)$ formé des $B'$ en question est cofinal dans $\Ss_{\pi}(B)$, et même $\pi$-cofinal.
\end{itemize}

Il faut prouver encore ce dernier point -- que $J$ (dont on sait qu'il est cofinal) est même $\pi$-cofinal, \textit{i.e.} stable par sup (dans $\Ss_{\pi}(B)$) des parties filtrantes $J' \subset J$ telles que $\mathrm{card}\, J' \leq \pi$. Mais cela résulte aussitôt du fait que $C$ étant stable par

\msp{page 83}
limites inductives filtrantes, $C^{\pi} = C \cap \Filt^{\pi}$ est stable par $\varinjlim$ filtrantes de cardinal $\leq \pi$, en vertu de la condition $(d)$ sur $\Filt^{\pi}$. Considérons maintenant le carré cocartésien
\[
\begin{tikzcd}
A' \arrow[r, hook, "i'"] \arrow[d, hook] & B' \arrow[d] & & \\
A \arrow[r, "i"] & C' = A \amalg_{A'} B' \arrow[r, "j"] & B \;.
\end{tikzcd}
\]
Comme $i' \in C^{\pi}$, on a $i \in \widetilde{C^{\pi}}$ (car $\widetilde{C^{\pi}}$ stable par cochangement de base). D'autre part, j'ai besoin du fait que $j$ soit dans $C$, et je ne sais pas même qu'il est un monomorphisme ! Aussi il faut repasser aux notations indicielles
\[
\begin{tikzcd}
A_i \arrow[r, hook, "\alpha_i"] \arrow[d, hook] & B_i \arrow[d] & \\
A \arrow[r, hook, "\beta_i"] & C_i \arrow[r, "\gamma_i"] & B \;.
\end{tikzcd}
\]
Comme la formation des sommes amalgamées commute aux limites inductives filtrantes, on voit que $\varinjlim C_i \simeq A \amalg_{\varinjlim A_i} \varinjlim B_i = A \amalg_A B = B$, donc la limite inductive des $\gamma_i$ est un isomorphisme (essentiellement $\mathrm{id}_B$). D'autre part, le même argument montre que le système inductif des $C_i$ $(i \in J)$ dans $\M$ est $\pi$-adapté, donc aussi celui des $\gamma_i$\footnote[28]{\textit{Sauf} qu'on ne sait pas que les $C_i$ soient dans $\Filt^{\pi}$ -- en fait, ils ne seront pas en général, puisqu'on ne suppose pas que $A, B \in \Filt^{\pi}$.}
\msp{page 84}
dans $\underline{\Fl}(\M)$. Comme $\varinjlim \gamma_i$ est un isomorphisme, il est dans $C$. Comme $C$ est $\pi$-accessible, il s'ensuit que l'ensemble $J'$ formé des $i \in J$ tels que $\gamma_i \in C$ est lui aussi cofinal\footnote[29]{\textit{Canulé, cf.} annotation précédente.}. Donc on trouve $i$ avec $\alpha_i \in C^{\pi}$, $\gamma_i \in C$, ce qui établit le lemme 6.4, et achève la démonstration du théorème 6.1.

\section*{7\quad Construction de triples de Quillen clos avec $\boldsymbol{C} = \mathbf{Mon} \subset \mathbf{Cof}_{\boldsymbol{W}}$}

Finalement, tous les beaux développements des paragraphes 5, 6 ne me donnent pas le “sweeping result” souhaité. Peut-être est-il faux -- peut-être la propriété d'engendrement
\printedpage{46}
% head: A. Grothendieck : Dérivateurs, Ch. XV, Édition des Universités de Montpellier II et Paris VII
“$^{\sim}$” par une petite partie (pour les ensembles accessibles, ou L-accessibles) est-elle beaucoup plus rarement satisfaite, que je ne l'aurais imaginé. Pour bien faire, il faudrait que je traite jusqu'au % resolved: the quoted raised wave denotes Q-closure; see corrections.csv, leaf 13.
\msp{page 85}
bout quelques exemples, p. ex. celui de $\Cof_W$, et de $W^{\mathrm{univ}}$ dans Cat, pour un localiseur fondamental plus ou moins général, et surtout dans le cas $W = \mathbf{W}_{\infty}$. Sûrement $\Cof_W$ et $W^{\mathrm{univ}}$ sont accessibles (il faudrait que je le vérifie pourtant), et si le “théorème” hypothétique du paragraphe 6 est faux pour eux, du moins en aurai-je le cœur net. Mais il reste pourtant la question : si $C = \widetilde{C}$ est stable par $\varinjlim$ filtrantes, et de la forme $\widetilde{C_0}$, $C_0$ petit, et si $W$ est L-accessible également, alors $C \cap W$ est-il aussi de la forme $\widetilde{TC_0}$, avec $TC_0$ petit ? Je viens encore d'essayer de le prouver, en supposant de plus $C \subset$ mono, et ne vois aucun moyen d'y parvenir. Il faudrait donc un contre-exemple pour dissiper ces perplexités ! Et le mieux serait dans Cat, tant qu'à faire.

Il surnage pourtant un résultat tangible positif, que je vais maintenant essayer de bien cerner et prouver :

\msp{page 86}
\textbf{Théorème 7.1.} \textit{Soit $\M$ une catégorie stable par petites $\varinjlim$ filtrantes, $M_0 \subset \Ob\,\M$ un sous-ensemble de} $\Ob\,\M$ \textit{générateur par monomorphismes stricts, plus précisément on suppose ceci :}

\begin{quote}
\itshape
Pour tout monomorphisme $i : A \hookrightarrow B$ dans $\M$ qui n'est pas un isomorphisme, il existe un monomorphisme $B' \overset{j}{\hookrightarrow} B$ ayant les propriétés suivantes :
\begin{itemize}
\item[\textup{a)}] $B' \in M_0$.
\item[\textup{b)}] $j$ ne se factorise pas par $i$.
\item[\textup{c)}] $A' = A \times_B B'$ existe dans $\M$, ainsi que $C = A \amalg_{A'} B'$, $A'$ est dans $M_0$, et dans le diagramme canonique
\[
\begin{tikzcd}
A' \arrow[r, hook, "i'"] \arrow[d, hook] & B' \arrow[d, hook] \arrow[dr, hook, dotted, "j"] & \\
A \arrow[r, hook, "\alpha"] \arrow[rr, hook, dotted, bend right=20, "i"'] & C \arrow[r, "\beta"] & B
\end{tikzcd}
\tag*{(7.1.1)}
\]
la flèche $\beta$ est un monomorphisme.
\end{itemize}
\end{quote}

(\textbf{N.B.} Ces conditions sont satisfaites si on suppose, en plus de la condition que $M_0$ soit générateur pour les monomorphismes, que $M_0$ contient, avec tout objet $B'$, et pour tout objet $A'$ qui est soit un quotient, soit un sous-objet de $B'$, un objet $A'_0$ isomorphe à $A'$ ; et si de plus $\M$ satisfait aux conditions d'exactitude suivantes :
\begin{itemize}
\item[1$^{\circ}$)] Toute flèche $f$ de $\M$ se factorise en $ip$, avec $i$ monomorphisme, $p$ épimorphisme, et :
\item[2$^{\circ}$)] Si $A_1$, $A_2$ sont deux sous-objets d'un objet $B$ de $\M$, alors l'intersection $A_0 = A_1 \times_B A_2$ existe, ainsi que la
\printedpage{47}
% head: A. Grothendieck : Dérivateurs, Ch. XV, Édition des Universités de Montpellier II et Paris VII
\msp{page 87}
somme amalgamée $C = A_1 \amalg_{A_0} A_2$, et le morphisme canonique $C \longrightarrow B$ est un monomorphisme.)
\end{itemize}

\textit{Supposons de plus que $\M$ satisfait la condition} \cercle{e}, p. 76, i.e. \textit{que toute $\varinjlim$ filtrante de monomorphismes est un monomorphisme.}

1$^{\circ}$) \textit{Soit} $\Mon \subset \Fl(\M)$ \textit{l'ensemble des monomorphismes, et} $\Mon_0$ \textit{le sous-ensemble de} $\Mon$ \textit{formé des $i \in$} $\Mon$ \textit{tels que source et but appartiennent à $M_0$. On a alors}
\[
\Mon \;\subset\; \widetilde{\Mon_0} \quad , \qquad \text{\textit{d'où}} \quad
\left\{
\begin{aligned}
\widetilde{\Mon} \quad &= \quad \widetilde{\Mon_0} \\
(\Mon)_{*} \quad &= \quad (\Mon_0)_{*} \;.\text{\footnotemark[30]}
\end{aligned}
\right.
\tag*{(7.1.2)}
\]
\footnotetext[30]{\textbf{N.B.} On a là un critère d'injectivité d'objets dans les catégories $\M/Y$ -- il suffit de tester la condition d'injectivité par des monomorphismes $i$ qui sont dans $\Mon_0$.}

2$^{\circ}$) \textit{Si $M_0$ est} petit \textit{et si les objets de $M_0$ sont accessibles} (p. ex. $\M$ satisfait \cercle{c}, p. 76), \textit{et si} $\Mon$ \textit{est stable par cochangement de base, alors il existe une factorisation fonctorielle de toute flèche $f$ de $\M$ en}
\[
f \;=\; p\, i \;,
\]
\textit{avec $i \in$} $\Mon$, \textit{$p \in$} $(\Mon)_{*}$. (\textbf{N.B.} Une flèche $p : X \longrightarrow Y$ de $\M$ est dans $(\Mon)_{*}$ si et seulement si elle fait de $X$ un objet \textit{injectif} de $\M/Y$ -- donc l'assertion de factorisation assure que dans les catégories induites $\M/Y$ “il y a assez d'injectifs”.) \textit{De plus, on a}
\[
\Mon \;=\; \widetilde{\Mon_0} \quad .
\tag*{(7.1.3)}
\]

3$^{\circ}$) \textit{Soit $W$ une partie de} $\Fl(\M)$ \textit{qui est} L-accessible (i.e. \textit{accessible, et stable par limites inductives filtrantes}). (\textit{On suppose à présent que $\M$ satisfait les conditions} (a), (b), (c), (d), (e) \textit{de la page 76, d'où une bonne théorie des parties L-accessibles de} $\Ob\,(\M)$ \textit{ou de} $\Fl(\M)$.) \textit{On suppose, plus précisément, que $W$ satisfait les conditions :}

\msp{page 88}
\[
\left\{
\begin{array}{@{\quad}l@{\quad}p{0.6\textwidth}}
\text{a)} & \textit{$W$ est L-accessible} (\textit{i.e.} stable par $\varinjlim$ filtrantes, et accessible). \\[1.5ex]
\text{b)} & \textit{Soient $\alpha$, $\beta$ composables, avec $\beta\alpha$ et $\alpha$ dans $W$, alors $\beta \in W$.} (\textbf{N.B.} Cette condition sera satisfaite si $W$ est un localiseur modérément saturé, et c'est à cette situation que nous allons appliquer notre énoncé.) \\[1.5ex]
\text{c)} & $\Mon \cap W$ \textit{est stable par cochangement de base.}
\end{array}
\right.
\tag*{(7.1.4)}
\]
\textit{et que $\M$ satisfait les conditions supplémentaires}
\[
\left\{
\begin{array}{@{\quad}l@{\quad}p{0.6\textwidth}}
\text{e$'$)} & \textit{Les $\varinjlim$ filtrantes commutent au changement de base.} \\[1.5ex]
\text{f)} & \textit{La condition d'exactitude} (7.2.1) \textit{ci-dessous.}
\end{array}
\right.
\tag*{(7.1.5)}
\]
\textit{Soit $\pi$ un cardinal tel que $W$ soit $\pi$-accessible, et $M_0 \subset$} $\Filt^{\pi}$ (\textit{donc $M_0$ essentiellement petite}). \textit{Alors on a}
\[
W \cap \Mon \;\subset\; \widetilde{(W \cap \Mon)^{\pi}}
\tag*{(7.1.6)}
\]
\printedpage{48}
% head: A. Grothendieck : Dérivateurs, Ch. XV, Édition des Universités de Montpellier II et Paris VII
(\textit{où pour toute partie $C$ de} $\Fl(\M)$, \textit{on pose} $C^{\pi} = C \cap \Fl^{\pi}(\M)$), \textit{donc}
\[
\widetilde{W \cap \Mon} \;=\; \widetilde{(W \cap \Mon)^{\pi}} \;,
\tag*{(7.1.7)}
\]
\textit{donc $\widetilde{W \cap \Mon}$ est Q-engendré par une} petite \textit{partie de} $W \cap \Mon$.

\textbf{Corollaire 7.2.} \textit{Soit $\M$ satisfaisant} (a), (b), (c), (d), (e) \textit{et aux conditions supplémentaires :}
\[
\left\{
\begin{array}{@{\quad}l@{\quad}p{0.6\textwidth}}
\text{e$'$)} & \textit{Les $\varinjlim$ filtrantes commutent au changement de base.} \\[1.5ex]
\text{f)} & \textit{Si $A_1$, $A_2$ sont deux sous-objets d'un objet $B$ de $\M$, et $A_0 = A_1 \cap A_2$, alors $A_1 \amalg_{A_0} A_2 \hookrightarrow B$ est un monomorphisme}\,\footnotemark[31].
\end{array}
\right.
\tag*{(7.2.1)}
\]
\footnotetext[31]{Une autre façon de le dire : si $C = \sup_{\Ss(B)}(A_1, A_2)$, le carré cartésien
\[
\begin{tikzcd}[ampersand replacement=\&]
A_0 \arrow[r] \arrow[d] \& A_1 \arrow[d] \\
A_2 \arrow[r] \& C
\end{tikzcd}
\]
est aussi cocartésien.}

\textit{Soit $W$ une partie de} $\Fl(\M)$ \textit{satisfaisant les conditions suivantes.}
\[
\left\{
\begin{array}{@{\quad}l@{\quad}p{0.6\textwidth}}
\text{a)} & \textit{$W$ est un localiseur modérément saturé.} \\[1.5ex]
\text{b)} & \textit{$W$ est L-accessible,} i.e. \textit{il est stable par $\varinjlim$ filtrantes, et est accessible.} \\[1.5ex]
\text{c)} & $\Mon \cap W$ \textit{est stable par cochangement de base,} i.e. \textit{si on a un diagramme} cocartésien
\begin{center}
\begin{tikzcd}[ampersand replacement=\&]
A \arrow[r, hook, "i"] \arrow[d, "\alpha"'] \arrow[dr, phantom, "\scriptstyle\text{cocart.}"] \& B \arrow[d, "\beta"] \\
A' \arrow[r, "i'"] \& B'
\end{tikzcd}
\end{center}
\textit{avec $i$ monomorphisme appartenant à $W$, alors $i'$ est un monomorphisme appartenant à $W$.}\,\footnotemark[32]
\end{array}
\right.
\tag*{(7.2.3)}
\]
\footnotetext[32]{\textit{Cf.} plus bas pour une quatrième condition $(d)$ sur $W$, qui assure que $W$ s'insère dans un triple de Quillen.}

\msp{page 89}
\textit{Sous ces conditions, on a}
\[
\underbrace{\Mon}_{C} \cap W \;=\; \widetilde{W \cap \Mon} \quad , \qquad \text{i.e. } \Mon \cap W \text{ \textit{est Q-clos à gauche}} \;,
\tag*{(7.2.4)}
\]

\end{document}
